\documentclass[11pt, letterpaper]{article}
\usepackage[utf8]{inputenc}
\usepackage[margin=1in]{geometry}
\usepackage{amsmath, amssymb}
\usepackage{booktabs}
\usepackage{tabularx}
\usepackage{xcolor}
\usepackage{enumitem}
\usepackage[font=bf, labelfont=bf]{caption}
\usepackage[protrusion=true, expansion=false]{microtype}
\usepackage{hyperref}
\usepackage{tcolorbox}
\tcbuselibrary{skins, breakable}

\definecolor{darkblue}{rgb}{0.0, 0.0, 0.55}
\definecolor{forestgreen}{rgb}{0.13, 0.54, 0.13}
\definecolor{charcoal}{rgb}{0.2, 0.2, 0.2}
\definecolor{revgray}{rgb}{0.96, 0.96, 0.98}
\definecolor{revborder}{rgb}{0.75, 0.78, 0.85}

\hypersetup{
    colorlinks=true,
    linkcolor=darkblue,
    filecolor=darkblue,
    urlcolor=darkblue,
    citecolor=darkblue
}

\newtcolorbox{reviewercomment}[1][]{
    enhanced,
    breakable,
    colback=revgray,
    colframe=revborder,
    coltitle=black,
    arc=2mm,
    boxrule=0.8pt,
    left=12pt,
    right=12pt,
    top=9pt,
    bottom=9pt,
    fontupper=\itshape,
    before skip=10pt,
    after skip=10pt,
    #1
}

\setlist[itemize]{leftmargin=1.5em, itemsep=3pt, topsep=4pt}
\setlist[enumerate]{leftmargin=1.5em, itemsep=3pt, topsep=4pt}

\setlength{\parskip}{6pt}
\setlength{\parindent}{0pt}

\begin{document}

{\LARGE \textbf{Response to Reviewers and Editor}}\\[6pt]
\textbf{Manuscript Number:} POETICS-D-26-00490\\
\textbf{Title:} \textit{Stay, React, or Conform? Examining Social Influence Effects on Aesthetic Judgments}\\
\textbf{Authors:} Omar Lizardo \& Sara Skiles\\
\textbf{Journal:} \textit{Poetics}\\
\textbf{Editor-in-Chief:} Hannah Wohl\\
\textbf{Date:} \today

\vspace{10pt}
\hrule
\vspace{15pt}

\section*{Overview and Letter to the Editor}

Dear Dr.~Wohl,

Thank you very much for the opportunity to revise and resubmit our manuscript, \textit{``Stay, React, or Conform? Examining Social Influence Effects on Aesthetic Judgments''} (POETICS-D-26-00490), for consideration in \textit{Poetics}. We are deeply grateful to you and the two anonymous reviewers for the constructive, insightful, and encouraging feedback.

We have carefully considered all comments from the editorial decision and the review reports. In response, we have developed a thorough revision addressing the methodological, conceptual, and empirical points raised.

Most centrally, our revision: (1)~\textit{clarifies the causal warrant and empirical hook} by explicitly positioning the study as providing an experimental causal test of taste influence and distancing dynamics that prior sociological literature had demonstrated primarily cross-sectionally or observationally; (2)~\textit{adopts a Difference-in-Differences (DiD) analytic framework} that benchmarks treatment effects against the unexposed control condition across trials, purging the estimates of counterfactual temporal drift and disambiguating baseline measurement from the control group; (3)~\textit{streamlines the theoretical framing} into a linear progression organized around \textit{valence asymmetry}, wherein negative evaluations act as an authoritative veto while positive evaluations facilitate appreciation only conditionally; (4)~\textit{articulates a sociological mechanism for status consistency}, drawing on status crystallization theory to explain why congruent positions foster evaluative certainty and behavioral stability (``staying''), whereas status inconsistency breeds fluidity; and (5)~\textit{expands methodological transparency} by providing ex-post Minimum Detectable Effect (MDE) power diagnostics, survey sampling details, descriptive statistics, exact treatment wording, full regression model tables, and expanded Inverse Probability Weighting (IPW) covariate balance diagnostics in a comprehensive Supplementary Appendix.

Below, we provide a point-by-point response to all comments and suggestions from the Editor and both Reviewers.

\vspace{15pt}
\hrule
\vspace{15pt}

\section*{Response to the Editor (Dr.~Hannah Wohl)}

\subsection*{1. Format Choice: Research Note vs.~Full Article}

\begin{reviewercomment}
While you originally submitted the manuscript as a ``research note,'' I offer two possible paths forward. You can revise the manuscript as a research note, honing in on the empirical contribution, or you can flesh out the manuscript and submit it as a full article, which would involve expanding on the theoretical contribution.
\end{reviewercomment}

\textbf{Location in Revised Manuscript:} Throughout revised manuscript (\texttt{manuscript\_R1.tex})

\textbf{Response:}
We are grateful to the Editor for giving us this choice. Following the guidance of both reviewers, we have pursued a path that captures the best of both options. First, we have kept a concise, focused manuscript structure centered squarely on the core empirical findings, while deepening the theoretical framework surrounding status crystallization, valence asymmetry, and causal identification.

To maintain narrative clarity and momentum in the main text, all extensive technical models, survey instrument vignettes, ex-post MDE power analyses, full regression tables, and auxiliary sensitivity checks have been organized into a Supplementary Appendix.

\vspace{10pt}

\subsection*{2. Empirical Novelty and Warrant for Research Question}

\begin{reviewercomment}
Reviewer 2 suggested that the empirical novelty and warrant for the research question could be made clearer. I think you could be more explicit about what specifically in the question, ``How does access to the evaluations of others influence our own aesthetic evaluations of cultural objects?'' has not been answered and/or how the literature offers competing answers. As I read it, you are offering the first empirical test of what can be inferred from a synthesis of the literature, but has not yet been empirically proven as such. If this is the case, I would make that clearer and update the literature (per Reviewer 1's suggestion).
\end{reviewercomment}

\textbf{Location in Revised Manuscript:} Section 1 (Introduction), Section 1.1--1.2 (Theoretical Background)

\textbf{Response:}
We thank the Editor for suggesting this framing. We have revised the Introduction and Section 1 to center directly on this empirical warrant. There, we note that decades of sociological scholarship (from Bourdieu and Simmel to contemporary cultural sociology) suggest that people's tastes should be directly influenced by information about others' tastes. We now explicitly frame the contribution as providing an experimental causal test of how evaluations of others---varying in both valence (like vs.~dislike) and status congruence---directly alter, suppress, or reinforce individual aesthetic judgments in real time.

\vspace{10pt}

\subsection*{3. Streamlining the Theoretical Front-End}

\begin{reviewercomment}
Reviewer 2 also notes that the takeaway and utility for researchers could be clearer. I think the issue might be that, because you brought together theories from multiple literatures to construct multiple hypotheses, your findings then affirm/complicate/challenge these multiple theories from multiple literatures, leaving the contribution a bit muddled. Streamlining the front end might help with this.
\end{reviewercomment}

\textbf{Location in Revised Manuscript:} Section 1.1--1.2 (Theoretical Framework), Section 1.4 (Hypotheses), Section 4 (Discussion)

\textbf{Response:}
We agree with the Editor and Reviewer 2 that weaving together disparate theoretical threads (such as informational cascades, Bourdieusian distinction, and status inconsistency) created a meandering opening in the original submission. In response, we have thoroughly streamlined the theoretical front-end of the manuscript into three tightly integrated, logically cumulative pillars.

First, we distinguish between purely informational and status-based social influence, clarifying how exposure to others' evaluations operates in the abstract compared to when evaluations are tethered to explicit status cues. Second, we formalize the logic of \textit{valence asymmetry}, articulating why negative consensus acts as an authoritative suppressive veto while positive consensus facilitates aesthetic appreciation only conditionally. Third, we ground our treatment of subjective class and education in Lenski's theory of status crystallization, explaining why structurally congruent positions (\textit{high-high} or \textit{low-low}) foster normative certainty and behavioral stability (``staying''), whereas status inconsistency generates evaluative fluidity. This streamlined organization establishes a direct and coherent throughline from theoretical expectations to our empirical hypotheses and results.

\vspace{10pt}

\subsection*{4. Methodological Rigor, DiD, and Demand Characteristics}

\begin{reviewercomment}
Finally, a successful revision will involve responding to both reviewers' methodology concerns, and the results will need to remain valid or be sufficiently updated. Reviewer 1 posed various methodological questions and offered suggestions for analyses, such as using the difference-in-differences approach, and asked for more context regarding the methodology. Reviewer 2's point about whether the experimental design might have been obvious to participants was something I wondered too, and providing more methodological detail might help alleviate these concerns.
\end{reviewercomment}

\textbf{Location in Revised Manuscript:} Section 2.1--2.7 (Data and Methods), Section 3.1--3.4 (Results), Section 4 (Discussion)

\textbf{Response:}
We have comprehensively addressed the methodological concerns raised by both reviewers while ensuring our core findings remain empirically rigorous and properly contextualized. Most substantially, we followed Reviewer 1's recommendation to reformulate our repeated-measures linear mixed models into a formal \textit{Difference-in-Differences (DiD)} framework. This specification explicitly benchmarks each experimental arm against the unexposed no-information control condition across pre- and post-treatment trials, purging the estimates of counterfactual temporal drift and exposing crucial causal patterns that were previously masked.

In addition, we have resolved concerns regarding potential demand characteristics and respondent reactivity (raised by both the Editor and Reviewer 2) by introducing Section 2.4. There, we document the four procedural safeguards embedded in the study design: the substantial 15-to-20 minute intervening survey buffer, the naturalistic framing of aggregate peer feedback, the unexposed control benchmark, and empirical behavioral patterns---namely, high baseline stability and active oppositional reactance---that directly contradict compliance with experimenter expectations.

\newpage

\section*{Response to Reviewer \#1}

\subsection*{Point 1: Analytical Strategy \& Difference-in-Differences (DiD) Approach}

\begin{reviewercomment}
1. The analytical strategy may not be the most appropriate. To better isolate the treatment effects, the authors may consider applying a difference-in-differences approach, using the no-information condition as the control group and exploiting the pre-treatment measurement. Without such a comparison, the observed changes could potentially be attributed to repeated measurement rather than the information treatment itself. The authors refer to two ``baselines''---the no-information condition and the measurement conducted before treatment---which suggests that a difference-in-differences design may be feasible.
\end{reviewercomment}

\textbf{Location in Revised Manuscript:} Section 2.7 (Difference-in-Differences Analytic Strategy), Section 3.1, Section 3.2, Figures 3--4, Table 2, and Appendix Tables A.2--A.3

\textbf{Response:}
We want to warmly and wholeheartedly thank the reviewer for their exceptional perceptiveness. Pushing us to adopt the Difference-in-Differences (DiD) framework was by far the most consequential and transformative recommendation we received. It not only rescued the paper from a fundamental inferential trap that we had failed to recognize in our original submission, but it also uncovered striking causal patterns that were previously masked.

\subsubsection*{1. What the Original Analysis Was Doing vs.~What DiD Accomplishes}
In our original submission (\texttt{manuscript.tex}), we evaluated simple pre-to-post change scores within each experimental cell:
\begin{equation*}
H_0: \Delta_{\text{within}} = (\bar{Y}_{k, \text{Trial 2}} - \bar{Y}_{k, \text{Trial 1}}) = 0
\end{equation*}
This one-sample test simply asked whether respondents in condition $k$ changed their aesthetic rating between Trial 1 and Trial 2. As the reviewer incisively observed, this one-group pre-post design conflates the true causal treatment effect of social feedback with the \textit{mere exposure effect}---the natural upward drift in appreciation that occurs when respondents re-evaluate an artwork upon repeated viewing.

Following the reviewer's guidance, we formulated a formal \textit{Difference-in-Differences (DiD) linear mixed model}:
\begin{equation}
\text{Taste}_{it} = \beta_0 + \beta_1 \text{Trial}_{it} + \sum_{k=2}^{K} \beta_k \text{Cond}_{ik} + \sum_{k=2}^{K} \delta_k (\text{Trial}_{it} \times \text{Cond}_{ik}) + u_i + \epsilon_{it}
\label{eq:did_model_resp}
\end{equation}
where the unexposed \textit{Control Condition} (which received no information between trials) serves as the omitted reference category. Here, the interaction coefficients $\delta_k$ directly isolate the causal treatment effect relative to counterfactual temporal drift:
\begin{equation*}
\delta_k = (\bar{Y}_{k, \text{Trial 2}} - \bar{Y}_{k, \text{Trial 1}}) - (\bar{Y}_{\text{Control}, \text{Trial 2}} - \bar{Y}_{\text{Control}, \text{Trial 1}})
\end{equation*}

\subsubsection*{2. Diagnosing What Changed: The Trap of Exposure Drift}
When we examined the unexposed Control Condition in our primary analytic sample (adults aged $\ge 24$, $N = 1,981$), the reviewer's intuition was powerfully confirmed: respondents who were given zero information nevertheless drifted upward substantially upon a second viewing. In the status-free baseline sample ($N = 457$), Control Condition respondents exhibited an upward exposure drift of $+0.156$ points ($SE = 0.055, t = 2.86, p = 0.0044$). Crucially, this exposure drift is heavily stratified across social strata: postgraduates drift upward by $+0.375$ points ($SE = 0.137, p = 0.0064$) and Bachelor's degree holders by $+0.149$ points ($SE = 0.101, p = 0.138$), whereas respondents with some college or associate degrees ($+0.021, p = 0.841$) and high school graduates ($+0.053, p = 0.651$) exhibit virtually zero drift. Similarly, across status consistency quadrants, Middle Class college graduates in the control group drift upward by $+0.250$ points ($p = 0.025$), whereas non-college working-class respondents show no change ($+0.016, p = 0.871$).

This discovery clarifies why several apparent positive effects from the original submission lost statistical significance in the DiD framework. In our initial draft, we reported that respondents exposed to high-status endorsements exhibited raw upward shifts, which we interpreted as evidence of upward conformity. However, once benchmarked against the counterfactual control group---who drifted upward on their own without seeing any social information---the net DiD estimates are statistically indistinguishable from zero (for instance, Bachelor's degree holders exposed to High-Status Likes yield $\text{DiD} = -0.030, p = 0.809$). What originally appeared to be active social conformity was, in reality, almost entirely an artifact of unmanipulated exposure drift. Furthermore, because DiD benchmarks directly against the control group ($N = 157$ in the primary sample), standard errors appropriately incorporate the variance of the control condition ($\text{Var}(\text{DiD}) = \text{Var}(\Delta_k) + \text{Var}(\Delta_{\text{Control}})$), increasing from approximately $0.06$ to $0.11\text{--}0.15$.

\subsubsection*{3. What Was Gained: New, Statistically Significant Causal Discoveries}
While spurious upward drift was properly eliminated, netting out counterfactual exposure drift uncovered major causal mechanisms that were entirely obscured in the original submission.

First, the DiD framework revealed vivid evidence of \textit{status-based distancing and taste abandonment (Hypothesis 5)}. In our revised four-tier educational analysis, respondents with a postgraduate degree exposed to \textit{Low-Status Likes} exhibited an immediate evaluative collapse of nearly half a point ($\text{DiD} = -0.491, SE = 0.135, t = -3.64, p = 0.0003$, Cohen's $d = 0.33$)---the single largest negative treatment effect in the experiment. Because postgraduates naturally drift upward by $+0.375$ under mere exposure, discovering that working-class peers endorsed the artwork actively reverses their appreciation, demonstrating Bourdieusian boundary-policing in real time. A parallel pattern emerges in the status consistency matrix, where Middle Class college graduates exposed to Low-Status Likes drop sharply ($\text{DiD} = -0.296, SE = 0.144, p = 0.0391$).

Second, the analysis establishes \textit{valence asymmetry as an authoritative suppressive veto (Hypothesis 2b)}. Across educational strata, negative evaluations consistently override natural exposure appreciation. In the status-free baseline sample, Generalized Dislike exerts a significant downward drag relative to control drift ($\text{DiD} = -0.186, p = 0.0169$). Among Bachelor's degree holders, both High-Status Dislikes ($\text{DiD} = -0.271, p = 0.0296$) and Generalized Dislikes ($\text{DiD} = -0.308, p = 0.0455$) induce significant evaluative suppression. Among postgraduates, this veto power is pervasive across all negative cues, including High-Status Dislikes ($\text{DiD} = -0.266, p = 0.0298$), Low-Status Dislikes ($\text{DiD} = -0.345, p = 0.0051$), and Generalized Dislikes ($\text{DiD} = -0.442, p = 0.0035$). By contrast, positive endorsements produce modest increments that fail to reach statistical significance beyond control drift ($p \ge 0.55$). Negative consensus thus acts as a potent veto on appreciation, whereas positive feedback is largely redundant.

Third, the DiD estimates provide a decisive empirical rejection of both \textit{cultural goodwill (Hypothesis 4)} and \textit{same-status alignment (Hypothesis 3)}. Lower-educated respondents (high school or less) remain completely unresponsive to elite cues (High-Status Like: $\text{DiD} = +0.032, p = 0.777$; High-Status Dislike: $\text{DiD} = -0.140, p = 0.215$), demonstrating aesthetic autonomy rather than deference. Likewise, respondents across all educational tiers fail to positively align with same-status peers (Bachelor's degree holders with High-Status Likes: $\text{DiD} = -0.030, p = 0.809$; high school graduates with Low-Status Likes: $\text{DiD} = +0.051, p = 0.697$).

\subsubsection*{4. Primary Difference-in-Differences Treatment Effects Table}
The table below presents the Difference-in-Differences treatment effects across the four educational tiers in the primary analytic sample (Age $\ge 24$, $N = 1,981$), estimated via multinomial Inverse Probability Weighting (IPW) balancing baseline age, gender, race/ethnicity, and parental education (Appendix Table A.2):

\begin{table}[ht!]
\centering
\small
\caption*{Difference-in-Differences Treatment Effects Across Four Educational Tiers (Age $\ge 24$ Sample, IPW-Weighted)}
\setlength{\tabcolsep}{3pt}
\begin{tabular*}{\linewidth}{@{\extracolsep{\fill}}lcccc@{}}
\toprule
\textbf{Treatment Condition} & \shortstack[c]{\textbf{HS or Less} \\ ($N=468$)} & \shortstack[c]{\textbf{2-Yr / Some Col} \\ ($N=655$)} & \shortstack[c]{\textbf{Bachelor's} \\ ($N=577$)} & \shortstack[c]{\textbf{Postgraduate} \\ ($N=281$)} \\
\midrule
Generalized Like & $+0.284^{\dagger}$ (0.151) & $+0.086$ (0.173) & $-0.040$ (0.145) & $-0.232$ (0.149) \\
High-Status Like & $+0.032$ (0.113) & $+0.145$ (0.127) & $-0.030$ (0.123) & $-0.235^{\dagger}$ (0.120) \\
Low-Status Like & $+0.051$ (0.130) & $-0.113$ (0.141) & $+0.044$ (0.139) & $-0.491^{***}$ (0.135) \\
Generalized Dislike & $-0.108$ (0.147) & $+0.069$ (0.154) & $-0.308^{*}$ (0.154) & $-0.442^{**}$ (0.151) \\
High-Status Dislike & $-0.140$ (0.113) & $-0.080$ (0.126) & $-0.271^{*}$ (0.125) & $-0.266^{*}$ (0.122) \\
Low-Status Dislike & $+0.240^{\dagger}$ (0.134) & $+0.036$ (0.139) & $-0.106$ (0.136) & $-0.345^{**}$ (0.123) \\
\midrule
Control Exposure Drift & $+0.053$ (0.118) & $+0.021$ (0.103) & $+0.149$ (0.101) & $+0.375^{**}$ (0.137) \\
\bottomrule
\end{tabular*}
\vspace{2pt}
{\footnotesize $^{\dagger}p < 0.10, ^{*}p < 0.05, ^{**}p < 0.01, ^{***}p < 0.001$. Standard errors in parentheses. Estimates represent DiD treatment contrasts relative to the unexposed Control Condition within each tier. The complete 24-cell status consistency quadrant matrix is reported in Appendix Table A.3.}
\end{table}

\subsubsection*{5. Streamlined and Publication-Ready Visualizations}
Rather than presenting a cluttered 24-cell grid in the main text, our revised empirical architecture is anchored by three focused figures. Figure 3 illustrates baseline exposure drift and generalized influence ($N = 457$), highlighting the mere exposure effect and the suppressive veto of Generalized Dislike. Figure 4 displays DiD treatment effects across the four educational tiers ($N = 1,981$), directly evaluating Hypotheses 3, 4, and 5. Figure 5 presents Average Marginal Effects for discrete behavioral choices (Stay, Conform, React) across status consistency quadrants, directly testing Hypothesis 6. In addition, Table 2 in the main text provides a comprehensive summary of every hypothesis (H1--H6), theoretical mechanisms, empirical estimates, and statistical verdicts, while full regression tables are preserved in Appendix C (Tables A.2 and A.3).

\vspace{10pt}

\subsection*{Point 2: Disambiguating ``Baseline'' vs.~``Control Condition'' Terminology}

\begin{reviewercomment}
2. I also suggest avoiding the term ``baseline'' when referring to the no-information condition, because the study includes a pre-treatment measurement that may also be understood as a baseline. Using the same term for both is confusing. The authors could instead refer to the no-information condition as the control condition.
\end{reviewercomment}

\textbf{Location in Revised Manuscript:} Throughout manuscript (\texttt{manuscript\_R1.tex}, Section 2.3, Section 3.1, and Section 3.2)

\textbf{Response:}
We agree completely with the reviewer that using the term ``baseline'' to denote both the temporal pre-test and the unexposed experimental cell created unnecessary ambiguity. In response, we have systematically overhauled our terminology throughout the manuscript.

Specifically, the no-information experimental group is now strictly referred to as the \textit{Control Condition} (or the \textit{no-information control group}). The initial, unmanipulated evaluation at Trial 1 ($t_1$) is designated as the \textit{pre-treatment baseline measurement}, and temporal changes over repeated trials in the absence of peer feedback are explicitly termed \textit{exposure drift} (or the \textit{mere exposure effect}). To ensure total clarity for readers from the outset, we have added an explicit note in Section 2.3 of the revised manuscript:
\begin{quote}
\textit{``To prevent terminological confusion, we strictly distinguish between the pre-treatment baseline measurement (Trial 1, which captures initial, uninfluenced aesthetic preference) and the Control Condition (the experimental group that receives no peer information between trials, allowing us to observe pure temporal re-evaluation in the absence of treatment). We refer to changes over repeated trials in this unexposed group as exposure drift.''}
\end{quote}

\vspace{10pt}

\subsection*{Point 3: Survey Sampling Procedure, Inference Scope, and Descriptive Statistics Table}

\begin{reviewercomment}
3. Please provide more details about the survey sampling procedure. Because this is an experimental study, external validity is less central than internal validity. However, if the survey used probability sampling, this would strengthen the conclusions. This information is also important for assessing what the reported confidence intervals represent---for example, whether they reflect only experimental uncertainty within the sample or can support population-level inference. Relatedly, the authors should include a table presenting descriptive statistics for the sample.
\end{reviewercomment}

\textbf{Location in Revised Manuscript:} Section 2.1 (Sample Recruitment and Survey Administration), Section 2.2 (Inferential Scope and Interpretation of Confidence Intervals), Table 1, and Appendix Table G.1

\textbf{Response:}
We have substantially expanded Section 2.1 and Section 2.2 to provide comprehensive details regarding the survey sampling procedure, participant recruitment, quota filtering, the primary analytic sample restriction, and the inferential scope of our statistical estimates.

Data were collected between February 28 and March 6, 2012, via an online survey administered by Survey Sampling International (SSI), drawn from its verified, opt-in consumer research panel. Participants provided informed consent and were compensated \$3.92 for completing the instrument (funded by NSF Doctoral Dissertation Improvement Grant SES-1203426 and the Notre Dame Institute for Scholarship in Liberal Arts). A total of 3,782 panel members accessed the survey; of these, 85 were excluded for non-consent, 12 were blocked for being under 18 years of age, and 1,394 were screened out because demographic quota targets---balanced to 2010/2012 U.S. Census benchmarks for age, gender, race/ethnicity, and geographic region---had already been met. An additional 16 participants dropped out during initial screeners, yielding a completed sample of $N = 2,275$ respondents ($N = 2,253$ providing complete pre- and post-treatment aesthetic judgments).

To ensure that all respondents had completed the traditional college-going window and eliminate right-censoring in educational attainment among the 18--24 cohort (where only 20.3\% held a bachelor's degree), we restrict our primary analytic sample to respondents aged 24 and older ($N = 1,981$ complete cases; 3,962 repeated trial observations). Within this primary sample, college degree holders represent 43.3\% ($N = 858$). The study implemented a deliberate oversampling of college graduates relative to national benchmarks ($\sim 28\%$) to secure adequate statistical power across high-cultural-capital occupational cells (Cultural Specialists and Economic Specialists). For completeness, the full unrestricted sample ($N = 2,275$) is preserved as a sensitivity check in Appendix G (Tables G.1, G.2, and G.3), confirming that all empirical findings and hypothesis verdicts are identical across the entire age spectrum.

Regarding the inferential scope of our statistics, because the survey utilized quota-matched non-probability panel sampling rather than probability sampling with known selection probabilities, the reported standard errors, confidence intervals, and $p$-values reflect \textit{model-based and experimental design uncertainty} within the sample rather than design-based population parameter inference. Nonetheless, the demographic diversity across gender, age, race, and geographic region provides substantially greater heterogeneity and ecological validity than conventional university convenience samples.

To document the sample composition transparently, we have added Table 1 in Section 2.1, displaying the demographic characteristics of the primary analytic sample alongside U.S. Census benchmarks and across experimental treatment and control arms:

\begin{table}[ht!]
\centering
\small
\caption*{Table 1: Demographic Characteristics of Survey Sample (Age $\ge 24$) Compared to U.S. Population Benchmarks}
\setlength{\tabcolsep}{2.5pt}
\begin{tabular*}{\linewidth}{@{\extracolsep{\fill}}lcccccc@{}}
\toprule
\textbf{Demographic Variable} & \textbf{Sample $N$} & \textbf{Sample \%} & \textbf{U.S. Pop \%} & \shortstack[c]{\textbf{Treatment} \\ ($N=1,524$)} & \shortstack[c]{\textbf{Control} \\ ($N=157$)} & \shortstack[c]{\textbf{Taste Only} \\ ($N=300$)} \\
\midrule
\textbf{Gender} & & & & & & \\
\quad Women & 1,036 & 52.3\% & 51.5\% & 51.9\% & 52.9\% & 54.0\% \\
\quad Men & 945 & 47.7\% & 48.5\% & 48.1\% & 47.1\% & 46.0\% \\
\addlinespace[2pt]
\textbf{Race / Ethnicity} & & & & & & \\
\quad White & 1,284 & 64.8\% & 65.1\% & 65.0\% & 64.3\% & 64.3\% \\
\quad Hispanic / Latino & 306 & 15.4\% & 15.8\% & 15.7\% & 14.6\% & 14.3\% \\
\quad Black / African American & 250 & 12.6\% & 12.3\% & 12.5\% & 13.4\% & 13.0\% \\
\quad Asian / Pacific Islander & 90 & 4.5\% & 4.5\% & 4.3\% & 5.7\% & 5.0\% \\
\quad Other / Multiracial & 51 & 2.6\% & 2.3\% & 2.5\% & 1.9\% & 3.3\% \\
\addlinespace[2pt]
\textbf{Age Group} & & & & & & \\
\quad 25--34 & 407 & 20.5\% & 17.5\% & 20.7\% & 19.1\% & 20.7\% \\
\quad 35--44 & 400 & 20.2\% & 17.5\% & 20.4\% & 15.3\% & 21.7\% \\
\quad 45--54 & 437 & 22.1\% & 19.2\% & 22.0\% & 23.6\% & 21.3\% \\
\quad 55--64 & 366 & 18.5\% & 15.6\% & 18.7\% & 17.2\% & 18.0\% \\
\quad 65+ & 371 & 18.7\% & 17.2\% & 18.2\% & 24.8\% & 18.3\% \\
\addlinespace[2pt]
\textbf{Geographic Region} & & & & & & \\
\quad Northeast & 371 & 18.7\% & 17.9\% & 19.0\% & 17.8\% & 17.7\% \\
\quad South & 689 & 34.8\% & 37.2\% & 35.6\% & 37.6\% & 29.3\% \\
\quad Midwest & 445 & 22.5\% & 21.6\% & 21.5\% & 26.1\% & 25.7\% \\
\quad West & 476 & 24.0\% & 23.3\% & 24.0\% & 18.5\% & 27.3\% \\
\addlinespace[2pt]
\textbf{Education (Oversampled)} & & & & & & \\
\quad Less than High School & 66 & 3.3\% & 13.3\% & 3.6\% & 1.9\% & 2.7\% \\
\quad High School Diploma & 402 & 20.3\% & 30.4\% & 19.8\% & 22.3\% & 22.0\% \\
\quad Some College / AA & 655 & 33.1\% & 28.5\% & 34.1\% & 30.6\% & 29.3\% \\
\quad Bachelor's Degree & 577 & 29.1\% & 18.1\% & 28.6\% & 29.9\% & 31.3\% \\
\quad Graduate / Professional & 281 & 14.2\% & 9.6\% & 14.0\% & 15.3\% & 14.7\% \\
\midrule
\textbf{Total $N$} & \textbf{1,981} & \textbf{100.0\%} & & \textbf{1,524} & \textbf{157} & \textbf{300} \\
\bottomrule
\end{tabular*}
\vspace{2pt}
{\footnotesize \textit{Note:} Full unrestricted sample demographics for $N = 2,275$, including the 18--24 cohort, are reported in Appendix Table G.1.}
\end{table}

\vspace{10pt}

\subsection*{Point 4: Sample Sizes \& Ex-Post Minimum Detectable Effect (MDE) Analysis}

\begin{reviewercomment}
4. The sample sizes within conditions are relatively small. If an a priori power analysis was not conducted during the study design, I recommend reporting an ex post minimum detectable effect (MDE) analysis in the appendix. This would help readers assess which effect sizes the study was capable of detecting. The authors may find the following discussion useful: ``Why ex post power using estimated effect sizes is bad, but an ex post MDE is not''.
\end{reviewercomment}

\textbf{Location in Revised Manuscript:} Section 2.2 (Inferential Scope and Interpretation of Confidence Intervals), Appendix A (Statistical Power and Minimum Detectable Effect Analysis), Table A.1, and Figure A.1

\textbf{Response:}
We are deeply grateful to the reviewer for this thoughtful methodological recommendation and for directing us to the literature distinguishing flawed retrospective \textit{post-hoc} power from principled \textit{ex-post} Minimum Detectable Effect (MDE) analysis (Bloom 1995; Gelman \& Carlin 2014; Hoenig \& Heisey 2001; McKenzie 2012). Because our sample size was fixed by grant funding during the 2012 survey administration, an ex-post MDE analysis provides the appropriate framework to assess what effect sizes the study was capable of detecting, without conditioning on observed effect estimates or sample $p$-values.

In response, we have added a comprehensive ex-post MDE power analysis in Appendix A of the revised manuscript, supported by a dedicated simulation and calculation script (\path{R/mde_power_analysis.R}), a summary table (Table A.1), and a dual-panel figure (Figure A.1) illustrating theoretical power curves alongside achieved empirical thresholds. Crucially, our power calculations account for the repeated-measures Difference-in-Differences structure of the data. Pre- and post-treatment aesthetic judgments exhibit exceptionally high test-retest autocorrelation ($r = 0.8993$, with $\sigma_1 = 1.504$ and $\sigma_2 = 1.547$). This strong within-subject stability shrinks the standard deviation of change scores to $\sigma_{\Delta} = 0.688$, reducing the standard error of treatment comparisons by more than 54\% relative to an independent post-test design ($\sigma_{\Delta} / \sigma_1 = 0.457$).

Assuming a two-tailed $\alpha = 0.05$ and standard 80\% statistical power, the MDE diagnostics establish several clear benchmarks. For within-subject exposure drift, the design was powered to detect changes as small as Cohen's $d \ge 0.029$ (0.043 scale points) in the pooled primary analytic sample ($N = 1,981$) and $d \ge 0.102$ (0.154 scale points) in the unexposed Control Condition ($N_C = 157$), confirming robust power to detect the observed $+0.156$ control drift ($p = 0.0044$). For pooled DiD contrasts against the Control Condition, the MDE thresholds are $d \ge 0.117$ (0.176 points) for pooled high-status conditions ($N_T = 505$), $d \ge 0.130$ (0.195 points) for single-status conditions ($N_T = 258$), and $d \ge 0.146$ (0.220 points) for taste-only conditions ($N_T = 150$). Because all pooled DiD models fall comfortably below Cohen's benchmark for a ``small'' effect ($d = 0.20$), the study was well powered to detect modest cultural influence.

Within status consistency quadrants, the median MDE ranges from $d \ge 0.223$ (0.335 scale points) for non-college working-class respondents and $d \ge 0.230$ (0.346 scale points) for college-educated middle-class respondents, up to $d \ge 0.285$ (0.429 scale points) for non-college middle-class and $d \ge 0.354$ (0.532 scale points) for college-educated working-class respondents. This analysis provides valuable interpretive clarity: the pooled models and status-consistent subgroups were adequately powered to detect small-to-moderate shifts, whereas smaller status-inconsistent cells required moderate-to-large behavioral responses to achieve statistical significance. As we note in the revised manuscript, null findings in the smallest cells reflect wider estimation uncertainty rather than affirmative evidence of zero influence.

\begin{table}[ht!]
\centering
\footnotesize
\caption*{Appendix Table A.1: Ex-Post Minimum Detectable Effect (MDE) Analysis (Age $\ge 24$ Sample)}
\setlength{\tabcolsep}{2.5pt}
\begin{tabular*}{\linewidth}{@{\extracolsep{\fill}}lccccccc@{}}
\toprule
& & & & \multicolumn{2}{c}{\textbf{80\% Power}} & \multicolumn{2}{c}{\textbf{90\% Power}} \\
\cmidrule(lr){5-6} \cmidrule(lr){7-8}
\textbf{Model Specification / Level} & $N_T$ & $N_C$ & $SE$ & $\text{MDE}_{\text{raw}}$ & Cohen's $d$ & $\text{MDE}_{\text{raw}}$ & Cohen's $d$ \\
\midrule
\multicolumn{8}{l}{\textit{Panel A: Within-Subject Exposure Drift (Paired Comparison)}} \\[2pt]
Primary Sample Drift & 1,981 & --- & 0.015 & 0.043 & 0.029 & 0.050 & 0.033 \\
Control Condition Drift & 157 & --- & 0.055 & 0.154 & 0.102 & 0.178 & 0.118 \\
\addlinespace[3pt]
\multicolumn{8}{l}{\textit{Panel B: Pooled Difference-in-Differences Contrasts (vs.~Control $N_C=157$)}} \\[2pt]
Taste-Only Conditions & 150 & 157 & 0.079 & 0.220 & 0.146 & 0.255 & 0.169 \\
Single Status Conditions & 258 & 157 & 0.070 & 0.195 & 0.130 & 0.226 & 0.150 \\
Pooled High-Status Conditions & 505 & 157 & 0.063 & 0.176 & 0.117 & 0.204 & 0.135 \\
\addlinespace[3pt]
\multicolumn{8}{l}{\textit{Panel C: Subgroup DiD Contrasts Within Status Consistency Quadrants}} \\[2pt]
Working Class, No College & 51--181 & 55 & 0.120 & 0.335 & 0.223 & 0.388 & 0.258 \\
Middle Class, College & 48--160 & 50 & 0.123 & 0.346 & 0.230 & 0.400 & 0.266 \\
Middle Class, No College & 26--125 & 31 & 0.153 & 0.429 & 0.285 & 0.496 & 0.330 \\
Working Class, College & 15--61 & 21 & 0.190 & 0.532 & 0.354 & 0.616 & 0.409 \\
\bottomrule
\end{tabular*}
\end{table}

\vspace{10pt}

\subsection*{Point 5: Exact Treatment Wording \& Vignette Prompts in Appendix}

\begin{reviewercomment}
5. Please report the exact wording of the information treatments in the appendix. For example, readers need to know whether the wording indicated that the reference group had the same or a different social status from the participants.
\end{reviewercomment}

\textbf{Location in Revised Manuscript:} Appendix B (Survey Vignettes and Treatment Prompts)

\textbf{Response:}
We have added the complete verbatim survey instrument flow, prompt wording, and treatment vignettes to Appendix B of the revised manuscript, transcribed directly from the original Qualtrics protocol.

In particular, addressing the reviewer's query regarding whether the reference group was framed as possessing the \textit{same} or \textit{different} status relative to the respondent, we clarify that the experimental vignettes \textit{never described the reference group in relative terms} (i.e., participants were never told that others had ``the same status as you'' or ``higher status than you''). Instead, the survey provided objective, absolute descriptions of the fictional peer group's educational credentials and occupational positions. For instance, the Low Cultural Capital vignette described prior respondents who ``completed high school but did not go to college, and work in a food service job.'' The Economic Specialist vignette specified peers who ``completed a Bachelor's degree in Business but no additional degrees, and work at a financial investment firm,'' while Cultural Specialists were described as having ``completed several advanced degrees including a PhD, and work at a college or university.'' In the Taste-Only condition, no socioeconomic information was provided (``Other people who have taken this survey recently have said that they [LIKE / DISLIKE] the painting you saw'').

Immediately following the vignette, respondents completed an attribution probe: \textit{``Given only what you know about this group's education, occupation, and opinion about the painting, how similar do you think you are to the people in this group?''} and were asked which dimension was most influential in their assessment of similarity (education/occupation, opinion of the painting, or both). The complete text and response scales for Trial 1, the intervening demographic and musical modules, and Trial 2 are now fully documented in Appendix B.

\vspace{10pt}

\subsection*{Point 6: Disaggregating Education and Subjective Class: Reorganizing Results and Expanding to Four Educational Tiers}

\begin{reviewercomment}
6. Before moving to the analysis of combined status, please report the results for subjective social class and education separately. These results could be presented as additional subfigures or included in the appendix. The robustness check using an alternative categorization is useful, as this possibility came to mind immediately when reading the analysis. It would also be relevant to mention this robustness check when the combined-status variable is first described.
\end{reviewercomment}

\textbf{Location in Revised Manuscript:} Section 2.5 (Variables and Measurement), Section 3.2 (Status-Differentiated Social Influence: Testing Hypotheses 3--5), Figure 4, Section 3.3 (Status Inconsistency and Discrete Behavioral Choices: Testing Hypothesis 6), Figure 5, Appendix Table A.2, Appendix Table A.3, Appendix Table G.2, and Appendix F

\textbf{Response:}
We are exceptionally grateful to the reviewer for this suggestion. Disaggregating educational attainment and subjective social class prompted us to fundamentally reconsider how structural status location is operationalized across our empirical pipeline. This led to a major, substantive reorganization that has vastly improved the clarity, theoretical coherence, and empirical power of the manuscript.

When we separated the two dimensions, two crucial insights became immediately apparent. First, examining subjective social class alone yielded an attenuated, diffuse picture with limited explanatory leverage on its own. Second, educational attainment emerged as the primary, unambiguous structural axis governing cultural capital, symbolic boundary work, and aesthetic judgment. This directly aligns with our experimental design, where the vignettes explicitly differentiated alter status through educational credentials and occupations (doctoral cultural specialists, business baccalaureates, and high school food service workers). Stratifying by education provides direct construct congruence between the status of the respondent (ego) and the status of the evaluator (alter).

Guided by the reviewer's recommendation, we examined whether our design was constrained to a blunt binary classification (College vs.~No College). Given our primary adult sample ($N = 1,981$ individuals; 3,962 repeated observations) and the high test-retest precision of our repeated-measures Difference-in-Differences design ($r = 0.899, \sigma_{\Delta} = 0.688$), we realized that we possess ample statistical power to move beyond a binary split and disaggregate educational attainment into four granular tiers: (1) High School or Less ($N = 468$); (2) Two-Year Associate Degree or Some College ($N = 655$); (3) Four-Year Bachelor's Degree ($N = 577$); and (4) Postgraduate or Advanced Degree ($N = 281$).

Estimating DiD models across these four tiers via generalized multinomial Inverse Probability Weighting (IPW) revealed a striking, monotonic cultural capital gradient, presented in Section 3.2 (Figure 4 and Appendix Table A.2). First, \textit{Hypothesis 3 (Same-Status Alignment) is cleanly rejected}: high-status respondents do not align with high-status likes (Bachelor's: $\text{DiD} = -0.030, p = 0.81$; Postgraduates: $\text{DiD} = -0.235, p = 0.05$), and lower-status respondents do not align with low-status likes ($\text{DiD} = +0.051, p = 0.70$) or dislikes ($\text{DiD} = +0.240, p = 0.074$). Second, \textit{Hypothesis 4 (Cultural Goodwill) is decisively rejected}: respondents with a high school education or less exhibit zero deference toward elite consensus, displaying flat shifts under High-Status Likes ($\text{DiD} = +0.032, p = 0.78$) and High-Status Dislikes ($\text{DiD} = -0.140, p = 0.22$). Instead, their responsiveness is selectively populist, shifting positively under Generalized Likes ($\text{DiD} = +0.284, p = 0.061$). Third, \textit{Hypothesis 5 (Status-Based Distancing) is strongly supported}: among Postgraduates, learning that working-class peers liked the artwork triggers an immediate evaluative collapse of nearly half a point ($\text{DiD} = -0.491, SE = 0.135, t = -3.64, p = 0.0003$, Cohen's $d = 0.33$)---the single largest negative effect in the entire experiment, providing vivid causal evidence of taste abandonment. Postgraduates also exhibit severe negative sensitivity across all dislike cues ($-0.266$ to $-0.442, p < 0.03$). Finally, the intermediate associate degree and some college cohort ($N = 655$) displays complete evaluative inertia across all conditions ($p \ge 0.25$), acting as an evaluative buffer between populist facilitation and elite boundary policing.

Having established the educational gradient, we reserve subjective social class for its true theoretical home in Section 3.3: cross-cutting education to construct the four-quadrant status consistency typology (Working/No College, Working/College, Middle/No College, Middle/College). Grounded in Lenski's (1954) status crystallization framework, status consistency governs normative certainty versus evaluative fluidity, which we evaluate directly through discrete behavioral choices via multinomial logistic regression (Figure 5). Status consistency anchors evaluative stability: Middle Class college graduates are $+7.3\%$ more likely to stay with their initial judgment ($\text{AME} = +0.073, z = 5.47, p < 0.001$, stay rate $81.1\%$), whereas inconsistent actors display lower stability. Conversely, status-inconsistent actors are significantly more susceptible to external influence (Working Class with College: $+4.2\%$ conformity, $p = 0.0003$; Middle Class without College: $+4.3\%$ conformity, $p = 0.0002$), while high-status consistent actors are significantly less likely to conform ($\text{AME} = -0.033, p = 0.0021$). Oppositional reactance is concentrated among non-college workers ($12.0\%$), whereas institutional cultural capital inhibits reactance ($-3.0\%$ and $-4.0\%, p < 0.003$).

This reorganization also allowed us to remove the old, unwieldy 24-cell continuous interaction matrix (old Figure 5) from the main text, which had diluted statistical power and invited multiple-comparison concerns. By relocating the 24-cell matrix to Appendix Table A.3 and reporting separate models for Education alone and Subjective Class alone in Appendix Table G.2, the empirical narrative flows seamlessly from baseline influence (Section 3.1) to educational status tiers (Section 3.2) to discrete status consistency choices (Section 3.3).

\vspace{10pt}

\subsection*{Point 7: Observational Nature of Status Categories \& IPW Logic}

\begin{reviewercomment}
7. Although the information treatment was randomly assigned, membership in categories such as ``low--high'' and ``low--low'' is also determined by respondents' pre-existing education or subjective social class. For example, individuals with lower levels of education may be more likely to occupy a lower position than the reference group---another reason why the exact treatment wording should be provided. Consequently, respondents do not necessarily have the same probability of entering each combined-status category, and comparisons across these categories may partly capture pre-existing differences between respondents rather than only the causal effect of the experimental information. This issue is common in studies that provide respondents with information about their position in the social structure (e.g., Cruces et al., 2013). The authors acknowledge this problem and address it using inverse probability weighting (IPW), but it would be valuable to devote a few additional sentences to explaining the logic of this approach and its implications for causal interpretation. Such a discussion would be a valuable contribution to the field and would help readers understand the rationale for the analytical strategy.
\end{reviewercomment}

\textbf{Location in Revised Manuscript:} Section 2.6 (Inverse Probability Weighting), Figure 2

\textbf{Response:}
We are deeply grateful to the reviewer for raising this vital methodological distinction and for directing us to the insightful framing in Cruces et al.~(2013). In Section 2.6 of the revised manuscript, we have added a dedicated, pedagogical explanation of the hybrid observational-experimental setup, the formal mechanics of Inverse Probability Weighting (IPW), and its exact implications for causal inference.

The core challenge arises from structural confounding: while the experimental treatments (peer likes vs.~dislikes and peer status vignettes) were randomly assigned, respondents' own structural status locations (objective education $\times$ subjective class) are observational attributes. As a result, respondents in different status groups differ systematically on baseline demographic characteristics. For example, Middle Class college graduates are substantially older and disproportionately have college-educated parents compared to non-college Working Class respondents. Unadjusted comparisons across status groups risk conflating true status effects with baseline demographic sorting.

Following Rosenbaum and Rubin (1983), we address this through a two-step IPW procedure. First, we fit a multinomial logistic regression predicting each respondent's probability of belonging to their observed status consistency quadrant $S_i \in \{1, 2, 3, 4\}$ (or educational group) conditional on strictly exogenous pre-treatment covariates $\mathbf{X}_i$ (age, gender, race/ethnicity, and parental education), yielding generalized propensity scores $e_s(\mathbf{X}_i) = P(S_i = s \mid \mathbf{X}_i)$. Second, each respondent is assigned a balancing weight equal to the inverse of their estimated propensity score: $w_i = 1 / e_s(\mathbf{X}_i)$.

Intuitively, this weighting creates a synthetic \textit{pseudo-population} by up-weighting underrepresented demographic profiles within each status group (such as working-class individuals whose parents attended college) and down-weighting overrepresented profiles. This severs the statistical dependency between the baseline covariates $\mathbf{X}_i$ and status group membership. Under the standard assumptions of conditional exchangeability (no unmeasured confounders given $\mathbf{X}_i$) and positivity ($0 < e_s(\mathbf{X}_i) < 1$), the weighted pseudo-population satisfies covariate balance. Consequently, subsequent differences in evaluative drift, behavioral inertia (``staying''), and reactive distancing across status categories can be causally attributed to structural status location and consistency, rather than demographic sorting.

\vspace{10pt}

\subsection*{Point 8: Covariate Balance Diagnostics}

\begin{reviewercomment}
8. The covariate-balance analysis should be explained more clearly. Does it compare the no-information control condition with all the other categories combined? Balance could also be assessed separately across all experimental groups, with the complete results reported in the appendix.
\end{reviewercomment}

\textbf{Location in Revised Manuscript:} Section 2.1, Table 1, Section 2.6 (Inverse Probability Weighting), Figure 2

\textbf{Response:}
We have revised Section 2.1 and Section 2.6 to clarify the exact comparisons evaluated across our balance diagnostics. First, regarding the experimental arms, Table 1 demonstrates that random assignment successfully balanced demographic characteristics (gender, race/ethnicity, age, region, and education) across the treatment, control, and taste-only conditions, with formal tests confirming no baseline imbalances (all $p > 0.40$).

Second, regarding the observational status categories, Figure 2 specifically evaluates covariate balance across the four status consistency quadrants (Middle Class/College, Middle Class/No College, Working Class/College, Working Class/No College). There, we compare each quadrant against the pooled population before and after weighting using the \texttt{cobalt} package (Greifer 2022). Prior to weighting, raw standardized mean differences (SMDs) for age and parental education ranged between 0.35 and 0.50 standard deviations. After applying IPW weights, all absolute SMDs are compressed well below the conservative 0.10 threshold across all covariates, confirming high balance in the weighted sample.

\vspace{10pt}

\subsection*{Point 9: Figure Polish: Dashed Line \& Capitalization}

\begin{reviewercomment}
9. In Figure 2, please explain what the dashed line represents and capitalize the variable labels consistently.
\end{reviewercomment}

\textbf{Location in Revised Manuscript:} Figure 2, Figure 3, Figure 4, Figure 5, Appendix Figure A.1, Appendix Figure A.2, Appendix Figure A.3

\textbf{Response:}
We have systematically overhauled all figures across the manuscript and appendix. All legacy forest plots have been replaced with clean, publication-standard horizontal bar plots featuring color-matched 95\% error bars, color-coded by statistical significance ($p < 0.05$ in red, $p < 0.10$ in orange, and non-significant in grey).

The dashed vertical line is now explicitly defined in every figure caption and plotted at zero ($\text{DiD} = 0$), representing the null benchmark of the unexposed Control Condition (or the standard $\pm 0.10$ threshold for covariate balance in Figure 2). In addition, standardized condition ordering and consistent title capitalization have been applied uniformly across all figures.

\vspace{10pt}

\subsection*{Point 10: Clarification of Generalized Influence Findings}

\begin{reviewercomment}
10. In discussing generalized influence, the authors state that ``exposure to generalized negative evaluations largely suppresses the positive baseline drift across all groups.'' If I understand correctly, readers are expected to interpret the ``dislike (no status)'' effect in Figure 3. However, only one effect appears to be statistically significant, and it is in the opposite direction. A similar issue arises for ``like (no status).'' I therefore wonder whether this interpretation is based on an analysis of the full sample that ignores the subgroup partition but is not currently reported or explained. The authors should clarify which results support this conclusion.
\end{reviewercomment}

\textbf{Location in Revised Manuscript:} Section 3.1 (Baseline Exposure Drift and Generalized Social Influence), Equation 3, Figure 3, Table 2

\textbf{Response:}
We appreciate the reviewer identifying this ambiguity in our initial draft. The earlier narrative conflated pooled models with subgroup interactions, creating confusion regarding which sample supported the generalized influence claims.

To resolve this, we introduced a dedicated subsection---\textit{Section 3.1: Baseline Exposure Drift and Generalized Social Influence}---that evaluates generalized influence cleanly on the status-free subset of the primary analytic sample ($N = 457$, comprising the Control Condition $N = 157$, Generalized Like $N = 152$, and Generalized Dislike $N = 148$). This analysis reveals two key empirical facts. First, in support of \textit{Hypothesis 1 (Mere Exposure Drift)}, repeated viewing in the absence of external social cues produces statistically significant upward appreciation ($\beta_1 = +0.156, SE = 0.055, t = 2.86, p = 0.0044$). Second, the data reveal a powerful \textit{valence asymmetry}: while learning that anonymous others liked the artwork yields an increase in raw evaluations ($+0.203$), this change is modest relative to the counterfactual $+0.156$ control drift ($\delta_{\text{Like}} = +0.047, SE = 0.079, t = 0.59, p = 0.555$, failing to support Hypothesis 2a). In stark contrast, learning that anonymous others disliked the artwork completely eliminates exposure appreciation (raw change $-0.030$), producing a statistically significant negative DiD suppression relative to control drift ($\delta_{\text{Dislike}} = -0.186, SE = 0.078, t = -2.40, p = 0.0169$, supporting \textit{Hypothesis 2b}).

A direct contrast between positive and negative cues yields a significant swing of nearly a quarter point ($\Delta = +0.233, SE = 0.080, \chi^2 = 8.51, p = 0.0035$). As visualized in Figure 3, this status-free baseline model provides clear, unambiguous support for Hypotheses 1 and 2b, demonstrating that negative consensus exerts an authoritative veto over natural aesthetic familiarization.

\vspace{10pt}

\subsection*{Point 11: Role of Disaggregated Behavioral Analysis}

\begin{reviewercomment}
11. The contribution of the disaggregated analysis is unclear to me. The categorical analysis also reports average effects, as do all the regression models in the paper. If the authors intend to move beyond average effects, they would need to apply methods that explicitly examine distributional differences or treatment-effect heterogeneity. I recommend removing this analysis, which would also free up space to develop other aspects of the research note.
\end{reviewercomment}

\textbf{Location in Revised Manuscript:} Section 3.3 (Status Inconsistency and Discrete Behavioral Choices), Figure 5, Appendix E, Appendix Figure A.3

\textbf{Response:}
We thank the reviewer for this thoughtful reflection. We carefully weighed removing the behavioral analysis, but chose to retain it in streamlined form because it resolves a critical limitation of continuous mean models. Specifically, in continuous models, an average treatment effect near zero cannot distinguish between complete evaluative stability (everyone staying) versus polarized churn (equal numbers conforming and reacting in opposite directions). By categorizing individual trajectories into mutually exclusive discrete modes---\textit{Stay} ($\Delta = 0$), \textit{Conform} (shifting in the direction of the cue), and \textit{React} (shifting against the cue)---multinomial logit models directly evaluate distributional shifts via Average Marginal Effects (AMEs).

This approach reveals three crucial empirical patterns that continuous means completely obscure, directly testing \textit{Hypothesis 6 (Status Consistency)}. First, status consistency anchors evaluative stability: high-status consistent individuals (Middle Class / College) exhibit the highest stay rate ($81.1\%$) and are significantly more likely to stay ($\text{AME} = +0.073, SE = 0.013, z = 5.47, p < 0.001$), while low-status consistent actors (Working Class / No College) also show high inertia ($73.8\%$ stay rate). In contrast, status-inconsistent actors display lower stability (Middle Class / No College stay rate $71.2\%, \text{AME} = -0.026, p = 0.070$). Second, status-inconsistent individuals are significantly more susceptible to external influence: Working Class college graduates are $+4.2\%$ more likely to conform ($\text{AME} = +0.042, p = 0.0003$), and Middle Class non-college respondents are $+4.3\%$ more likely to conform ($\text{AME} = +0.043, p = 0.0002$), whereas high-status consistent respondents are significantly less likely to conform ($\text{AME} = -0.033, p = 0.0021$). Third, institutional cultural capital sharply inhibits oppositional reactance ($-3.0\%$ and $-4.0\%, p < 0.003$), whereas reactance is concentrated among non-college workers ($12.0\%$).

To address concerns regarding robustness, we also report a dedicated sensitivity analysis in Appendix E (Appendix Figure A.3) excluding the Taste-Only conditions. This confirms that the discrete behavioral distributions replicate with exceptional fidelity when only status-laden cues are evaluated.

\vspace{10pt}

\subsection*{Point 12: Full Regression Model Tables in Appendix}

\begin{reviewercomment}
12. Report tables with full models in the appendix.
\end{reviewercomment}

\textbf{Location in Revised Manuscript:} Table 1, Table 2, Appendix Table A.1, Appendix Table A.2, Appendix Table A.3, Appendix Table A.4, Appendix Table G.1, Appendix Table G.2, Appendix Table G.3

\textbf{Response:}
We have fully satisfied this request by incorporating complete regression tables and diagnostic metrics across the manuscript and Supplementary Appendix. In the main text, Table 1 reports descriptive statistics for the primary analytic sample ($N = 1,981$) against 2012 Census benchmarks across experimental conditions, and Table 2 provides a comprehensive summary of all hypotheses (H1--H6), theoretical mechanisms, empirical estimates, and statistical verdicts.

The Supplementary Appendix provides the full suite of econometric models: Table A.1 reports the complete ex-post Minimum Detectable Effect (MDE) power analysis; Table A.2 presents the full Difference-in-Differences regression models across the four educational tiers; Table A.3 details the complete 24-cell DiD matrix across status consistency quadrants; and Table A.4 provides a side-by-side comparison of unweighted and IPW-weighted pooled DiD models with random effect variance components. In addition, Appendix G provides the full unrestricted sample models ($N = 2,275$), including complete descriptive statistics (Table G.1), stepwise DiD models separating education and class (Table G.2), and the 24-cell status consistency matrix (Table G.3).

\vspace{10pt}

\subsection*{Point 13: Engagement with Recent Taste \& Influence Debates}

\begin{reviewercomment}
13. My only comment on the theoretical discussion is that the authors should engage more directly with the latest debates on taste and social influence, including more recent papers and book chapters by Lizardo and other scholars working in this area.
\end{reviewercomment}

\textbf{Location in Revised Manuscript:} Section 1 (Introduction), Section 1.1--1.2 (Theoretical Background), Section 4 (Discussion), References

\textbf{Response:}
We have thoroughly updated our theoretical framing and bibliography to engage recent scholarship on cultural cognition, dual-process habitus, aesthetic signaling, and symbolic boundaries (e.g., Lizardo 2016; Lizardo et al.~on the cognitive foundations of cultural tastes; Wohl 2015; Vanzella 2022). This contemporary literature reinforces our core argument: aesthetic evaluations do not merely reflect internal hedonic reactions, but serve as cognitive anchors that interface with external social cues to signal group belonging, status distinction, and structural boundary policing.

\newpage

\section*{Response to Reviewer \#2}

\subsection*{Point 1: Research Question Rationale \& Causal Hook}

\begin{reviewercomment}
In section 1.1, the paper reviews past research that shows that people tend to share tastes with similar others and tend to have distaste for the tastes of dissimilar others. (I would agree that this is the state of the literature.) The authors then state the research question: ``How does access to the evaluations of others influence our own aesthetic evaluations of cultural objects?'' It's not clear why this research question needs to be asked, given the literature that just got reviewed and the authors' characterization of that literature. However, I think there is room to ask whether these relationships have been rigorously demonstrated with data that are designed to assess causal relationships. I think it might be the case that past work is observational and cross-sectional. Perhaps that should be the hook that highlights the novelty and contribution of the findings?
\end{reviewercomment}

\textbf{Location in Revised Manuscript:} Section 1 (Introduction), Section 1.1--1.2 (Theoretical Background)

\textbf{Response:}
We are deeply grateful to Reviewer 2 for this crucial insight. We have restructured the entire opening of the paper around this exact methodological and theoretical hook. We now explicitly emphasize that while observational studies establish homophilous taste patterns, they cannot determine whether shared tastes reflect prior sorting, shared structural environments, or active interpersonal influence and reactive distancing. Our study addresses this gap by providing an experimental causal test of these micro-interactional dynamics in real time.

\vspace{10pt}

\subsection*{Point 2: Streamlining Section 1}

\begin{reviewercomment}
The third paragraph of section 1.1 reviews work showing that high status groups' tastes tend to induce conformity in others' tastes and low status groups' tastes tend to induce distancing. The fourth paragraph brings the status of the self back in, which gets us back to the starting point, where taste conformity results from social similarity between groups and taste distancing results from social dissimilarity. It is not clear to me that the self was ever not in the analysis. The last paragraph of section 1.1 adds in a new consideration: the distinction between objective and subjective class status. For me, this distinction is itself distinct from the material that has come before, which only deals with objective social similarity and dissimilarity. I think section 1.1 has to be rewritten to be more focused and streamlined on the concepts central to the data analysis. Right now it is meandering.
\end{reviewercomment}

\textbf{Location in Revised Manuscript:} Section 1.1--1.3 (Theoretical Background)

\textbf{Response:}
We agree with the reviewer that weaving together multiple theoretical threads produced conceptual loops in the original draft. We have completely rewritten Section 1 to follow a linear, cumulative progression across three central concepts.

We begin with the distinction between informational and status-based social influence, explaining how cues from others act as heuristics under aesthetic ambiguity while remaining fundamentally structured by valence asymmetry and alter status. Next, we address relational status matching, examining how influence is governed by the structural congruence between ego's social location and alter's status (same-status alignment versus cross-status reactance). Finally, we introduce status consistency, demonstrating how structural alignment between objective credentials and subjective class anchors cognitive certainty and evaluative inertia (``staying''), whereas status inconsistency generates evaluative fluidity.

\vspace{10pt}

\subsection*{Point 3: Demand Characteristics and Respondent Reactivity}

\begin{reviewercomment}
My understanding of the data being analyzed is that they are responses to a question on liking of the painting, compared to a repeated question on liking that followed the provision of information about who else liked/disliked the painting, with the occupation and education of those others being varied. I'm wondering whether this design made respondents aware that a shift in their preference was being measured, and how that might influence their responses. Where any steps taken to minimize demand awareness or respondent reactivity?
\end{reviewercomment}

\textbf{Location in Revised Manuscript:} Section 2.4 (Mitigation of Experimenter Demand Characteristics and Respondent Reactivity)

\textbf{Response:}
We thank Reviewer 2 for raising this important methodological consideration. In response, we have added a dedicated subsection---Section 2.4---detailing four procedural and design safeguards that minimized demand awareness and respondent reactivity.

First, pre- and post-treatment evaluations were separated by an extensive 15-to-20 minute intervening survey buffer comprising detailed employment questions, parental demographics, and a comprehensive 20-genre musical taste battery, preventing mechanical recall of initial ratings. Second, aggregate peer feedback was presented naturalistically as observational survey summaries followed by subjective similarity and attribution probes, keeping the apparent purpose focused on social perception rather than persuasion. Third, benchmarking all treatment conditions against the unexposed Control Condition rigorously separates treatment effects from any potential repeated-testing reactivity. Fourth, the empirical behavioral distributions directly contradict passive compliance with experimenter demand: over 74\% of respondents exhibited complete evaluative stability (``staying''), and significant oppositional reactance emerged among discordant status groups.

\vspace{10pt}

\subsection*{Point 4: Framing of Valence Findings: Valence Asymmetry}

\begin{reviewercomment}
In section 4, the authors write, ``While previous research suggests that information about others' choices acts as a heuristic to resolve uncertainty (Strang and Macy, 2001; Salganik et al., 2006), we find that valence matters.'' Isn't it true that valence mattered in past research as well, insofar as others' likes vs dislikes mattered? Instead, isn't the novel finding here is that there is an asymmetry to the valence?
\end{reviewercomment}

\textbf{Location in Revised Manuscript:} Abstract, Section 1.2, Section 3.1, Section 4.1

\textbf{Response:}
We thank Reviewer 2 for this sharp and constructive correction. We have revised our framing throughout the manuscript to focus specifically on \textit{valence asymmetry} rather than suggesting that prior literature ignored valence altogether. We now explicitly highlight that while positive peer feedback produces selective, status-contingent shifts, negative peer feedback acts as an authoritative asymmetric veto that broadly halts positive exposure drift across groups.

\vspace{10pt}

\subsection*{Point 5: Theoretical Mechanism for Status Inconsistency as a Driver of Evaluative Fluidity}

\begin{reviewercomment}
Although the discussion very briefly describes status inconsistency as a driver of evaluative fluidity, that finding is not sufficiently integrated with the hypotheses and there is no clear mechanism attributed to it. The authors should formally test whether status inconsistency moderates the experimental conditions. Is the paper also meant to demonstrate that there is a general propensity to change that is distinct from conformity to particular status groups? If so, there needs to be more explanation in the discussion to describe that and to explain the significance of that.
\end{reviewercomment}

\textbf{Location in Revised Manuscript:} Section 1.2, Section 1.4 (Hypothesis 6), Section 3.3, Section 4.3

\textbf{Response:}
We have fully articulated the theoretical mechanism grounded in Lenski's (1954) status crystallization theory and contemporary cultural sociology. Actors occupying congruent structural locations (\textit{high-high} or \textit{low-low}) are embedded within coherent social milieus with unambiguous normative expectations, generating high subjective certainty and behavioral stability (``staying''). In contrast, when actors possess college degrees while identifying as working class, or claim middle-class identity without institutional credentials, they experience cross-pressured reference norms. This structural ambiguity lowers evaluative inertia and heightens susceptibility to external cues, leading to elevated conformity and reactivity.

We formally evaluate this mechanism in Section 3.3 via Average Marginal Effects from multinomial logistic regression models. The empirical results confirm that status-inconsistent groups are significantly more likely to conform (Working Class with College: $+4.2\%, p = 0.0003$; Middle Class without College: $+4.3\%, p = 0.0002$) and display significantly lower evaluative stability than consistent groups.

\vspace{10pt}

\subsection*{Point 6: Clarification of ``Stay'' / Inertia Across Consistent Groups}

\begin{reviewercomment}
This sentence needs clarification: ``For instance, we see that status consistent individuals, and particularly high-status consistent individuals have the highest propensity to stay with their initial judgment of taste (Middle Class/College: 79.0\%; Working Class/No College: 74.4\%), a classic behavioral indicator of high status (Correll and Ridgeway, 2003).'' The lowest status individuals have a much higher propensity to stay than the two inconsistent groups who both have higher status. What does this imply about the link between high status and staying with a taste judgment? And it amplifies the need for a mechanism to explain the link between status inconsistency and evaluative fluidity.
\end{reviewercomment}

\textbf{Location in Revised Manuscript:} Section 3.3, Section 4.3 (Discussion)

\textbf{Response:}
Reviewer 2 makes an exceptional point. We have corrected the narrative to clarify that evaluative stability (``staying'') is not a monotonic function of high status alone, but rather an outcome of \textit{structural status consistency and normative certainty}. Both high-consistent individuals (Middle Class / College: 81.1\% stay rate) and low-consistent individuals (Working Class / No College: 73.8\% stay rate) occupy crystallized positions within the social structure, providing high evaluative certainty. In contrast, status-inconsistent individuals (Middle Class / No College: 71.2\% stay rate, $\text{AME} = -0.026, p = 0.070$; Working Class / College: elevated conformity $+4.2\%, p = 0.0003$) experience structural cross-pressures, resulting in lower inertia and greater evaluative fluidity.

\vspace{10pt}

\subsection*{Point 7: Clarity and Utility for Researchers}

\begin{reviewercomment}
For me, one of the virtues of a research note is that it presents clear findings that other researchers can rely on for having established an empirical set of facts. If find that this research note is lacking that clarity, so I'm not sure how readily it will be useful to other researchers.
\end{reviewercomment}

\textbf{Location in Revised Manuscript:} Abstract, Section 3.4, Table 2, Section 4.1 (Discussion)

\textbf{Response:}
To ensure maximum clarity and utility for other researchers, we have synthesized the study's core empirical findings into \textit{Table 2 (Summary of Hypotheses, Expectations, and Empirical Verdicts)} in Section 3.4. This table presents an unambiguous summary of the theoretical expectations and empirical verdicts across all six hypotheses.

First, regarding counterfactual exposure drift, liking drifts upward significantly in the absence of external cues ($\beta_1 = +0.156, p = 0.0044$), driven primarily by postgraduates ($\beta_1 = +0.375, p = 0.0064$) and bachelor's degree holders ($+0.149$), while non-college respondents exhibit no drift ($+0.053, p = 0.651$). Second, demonstrating valence asymmetry as an asymmetric veto, negative evaluations actively suppress positive exposure drift ($\delta_{\text{Dislike}} = -0.186, p = 0.0169$; High-Status Dislike among BA holders $\text{DiD} = -0.271, p = 0.0296$; High-Status Dislike among postgraduates $\text{DiD} = -0.266, p = 0.0298$), whereas positive evaluations fail to reach significance beyond exposure drift ($\delta_{\text{Like}} = +0.047, p = 0.555$). Third, demonstrating Bourdieusian cross-status reactance and taste abandonment, high-status postgraduates devalue and abandon an artwork when informed that working-class peers liked it ($\text{DiD} = -0.491, p = 0.0003$, Cohen's $d = 0.33$), with a parallel drop among Middle Class college graduates ($\text{DiD} = -0.296, p = 0.0391$). Finally, evaluative stability is anchored by structural status consistency (stay rates of 81.1\% and 73.8\%), whereas status inconsistency significantly elevates conformity ($+4.2\%$ and $+4.3\%, p < 0.001$).

\end{document}
